\begin{document}

\title{从零开始使用 MetaBlog}
\author{MetaBlog 示例}{}{example@example.com}

\begin{abstract}
这篇文章介绍一个 MetaBlog 站点所需的最小文件结构，以及构建和预览站点时最常用的命令。
\end{abstract}

\begin{keywords}
MetaBlog, 网站构建, 文章元数据
\end{keywords}


\section{项目地址}

\href{https://github.com/MetaronWang/MetaBlog}{\textbf{MetaBlog}, An awesome TeX-based blog website cli}

\section{目录结构}

一个 MetaBlog 站点就是一个普通文件夹。最重要的文件包括：

\begin{description}
  \item[\texttt{data/config.toml}] 网站标题、logo、icon 和分页配置。
  \item[\texttt{data/articles.toml}] 需要参与构建的文章列表。
  \item[\texttt{data/about_page/main.tex}] 关于页面的 LaTeX 源文档。
  \item[\texttt{articles/}] 文章源文件目录。
  \item[\texttt{asset/}] 站点级资源，例如 logo 和 favicon。
\end{description}

只有注册在 \texttt{data/articles.toml} 中的文章才会出现在生成的网站中。\texttt{articles/} 下存在某个文件夹，并不意味着它一定会被构建。

\Section{初始化站点}

通过golang运行源代码，或使用我们发布的二进制程序，可以在任意空目录初始化网站

\begin{lstlisting}[language=PowerShell]
go run ./cmd/metablog site init -root example -title "MetaBlog Example"
# 使用release的二进制包：
./metablog_1.0.2_windows_amd64.exe site init -root example -title "MetaBlog Example"
\end{lstlisting}

\texttt{site init} 默认会下载字体文件，并在最后检测 Python、LaTeXML 和 PDF 转换器环境。它只检测并提示安装方式，不会自动安装或修改系统环境。

如果只想先生成目录和配置，可以跳过字体下载和环境检测：

\begin{lstlisting}[language=PowerShell]
go run ./cmd/metablog site init -root example -title "MetaBlog Example" -skip-fonts -skip-env-check
# 使用release的二进制包：
./metablog_1.0.2_windows_amd64.exe site init -root example -title "MetaBlog Example" -skip-fonts -skip-env-check
\end{lstlisting}

\section{构建和预览}

\href{https://github.com/MetaronWang/MetaBlog}{Github仓库}中的 \texttt{example/}文件夹 是一个完整示例站点，
站点名为 \texttt{MetaBlog Example}，包含关于页、自定义组件和多篇展示文章。可以直接从仓库根目录构建或预览：

从仓库根目录构建这个示例：

\begin{lstlisting}[language=PowerShell]
go run ./cmd/metablog site build -root example -out example/out
# 使用release的二进制包：
./metablog_1.0.2_windows_amd64.exe site build -root example -out example/out
\end{lstlisting}

输出目录是 \texttt{example/out}。如果需要实时预览，可以使用：

\begin{lstlisting}[language=PowerShell]
go run ./cmd/metablog site serve -root example -out example/out -watch
# 使用release的二进制包：
./metablog_1.0.2_windows_amd64.exe site serve -root example -out example/out -watch
\end{lstlisting}

启用 \texttt{-watch} 后，已注册文章目录、关于页面、站点元数据、站点资源和自定义组件都会被监听。文档变化后，MetaBlog 会增量重建对应页面。

\section{文章元数据}

每篇文章都通过 TOML 配置：

\begin{lstlisting}[language=TOML]
[[articles]]
title = "从零开始使用 MetaBlog"
description = "这篇文章介绍一个 MetaBlog 站点的基本结构。"
author = "MetaBlog 示例"
date = "2026-05-16"
category = ["指南", "快速开始"]
tags = ["MetaBlog", "CLI", "网站构建"]
folder = "articles/getting-started-with-metablog"
main_fig = "figs/site-overview.svg"
main_file = "main.tex"
slug = "getting-started-with-metablog"
\end{lstlisting}

\texttt{description} 会显示在首页和文章列表页中。\texttt{category} 是多级分类，\texttt{tags} 中的每个标签都会生成对应的标签列表页。

\section{分页}

这个示例把 \texttt{home_page_size} 设置为 3，并提供 4 篇文章，因此首页可以直接展示分页效果。

\end{document}
